\DocumentMetadata{
  pdfversion=2.0,
  pdfstandard=ua-2,
  lang=en-US,
  tagging=on,
  tagging-setup={math/setup=mathml-SE,tabsorder=structure}
}
\documentclass[11pt,letterpaper]{article}
\usepackage[margin=0.68in,includefoot]{geometry}
\usepackage{newcomputermodern}
\usepackage{amsmath,amscd,mathtools}
\usepackage{tikz,adjustbox}
\providecommand{\square}{\mdlgwhtsquare}
\usepackage[hidelinks]{hyperref}
\hypersetup{
  pdftitle={Complex Analysis PhD Exam, January 2014},
  pdfauthor={University of Florida Department of Mathematics},
  pdfsubject={Graduate mathematics examination},
  pdfdisplaydoctitle=true,
  bookmarksopen=true
}
\setcounter{secnumdepth}{0}
\setlength{\parindent}{0pt}
\setlength{\parskip}{0.28em}
\setlength{\emergencystretch}{3em}
\setlength{\leftmargini}{2.15em}
\setlength{\leftmarginii}{2.65em}
\setlength{\leftmarginiii}{3em}
\pagestyle{empty}
\raggedbottom
\allowdisplaybreaks
\NewTaggingSocketPlug{block/list/label}{examproblem}{%
  \tagstructbegin{tag=Lbl}\tagstructend%
  \tagstructbegin{tag=\LBody}%
  \tagstructbegin{tag=\ProblemHeadingTag,title-o={\ProblemHeadingText}}%
  \tagmcbegin{tag=\ProblemHeadingTag,actualtext={\ProblemHeadingText}}%
  #2%
  \tagmcend%
  \tagstructend%
}
\newenvironment{examproblems}[2]{%
  \def\ProblemHeadingTag{#2}%
  \AssignTaggingSocketPlug{block/list/label}{examproblem}%
  \begin{enumerate}%
  \setcounter{enumi}{#1}%
  \renewcommand{\labelenumi}{\textbf{\theenumi.}}%
  \setlength{\topsep}{0.35em}%
  \setlength{\partopsep}{0pt}%
  \setlength{\itemsep}{0.48em}%
  \setlength{\parsep}{0.18em}%
}{\end{enumerate}}
\newenvironment{examsubparts}{%
  \AssignTaggingSocketPlug{block/list/label}{default}%
  \begin{enumerate}%
  \setlength{\topsep}{0.18em}%
  \setlength{\partopsep}{0pt}%
  \setlength{\itemsep}{0.16em}%
  \setlength{\parsep}{0.1em}%
}{\end{enumerate}}
\newenvironment{examinstructions}{%
  \par\begingroup\itshape
}{%
  \par\endgroup\vspace{0.18em}
}
\makeatletter
\renewcommand\subsection{\@startsection{subsection}{2}{0pt}%
  {-1.25ex plus -.25ex minus -.1ex}%
  {0.55ex plus .1ex}%
  {\normalfont\large\bfseries}}
\renewcommand{\@maketitle}{%
  \newpage\null\vskip -1em%
  {\footnotesize\bfseries
    \href{https://www.ufl.edu/}{University of Florida}, \href{https://math.ufl.edu/}{Department of Mathematics}\par}%
  \vskip -0.5em%
  \tagstructbegin{tag=H1,title={Complex Analysis PhD Exam, January 2014}}%
  \tagpdfparaOff%
  \tagmcbegin{tag=H1}%
    {\LARGE\bfseries\href{https://gma.math.ufl.edu/exams/complex-analysis-phd/}{Complex Analysis PhD Exam}, January 2014\par}%
  \tagmcend%
  \tagpdfparaOn%
  \tagstructend%
  \vskip 0.25em%
  \tagmcbegin{artifact}%
    \hrule height 1pt%
  \tagmcend%
  \vskip 1em%
}
\makeatother
\title{Complex Analysis PhD Exam, January 2014}
\author{}
\date{}
\begin{document}
\maketitle
\thispagestyle{empty}
\begin{examinstructions}
Do any 9 problems. The notation \(f \in H(\Omega)\) means~\(f\) is holomorphic in the domain~\(\Omega\).
\end{examinstructions}
\begin{examproblems}{0}{H2}
\def\ProblemHeadingText{Problem 1.}
\item
Prove the fundamental theorem of algebra.
\def\ProblemHeadingText{Problem 2.}
\item
Let~\(\Omega\) be a domain on the complex plane and~\(M\) be a finite positive constant. Suppose~\(F\) is a family of holomorphic functions on~\(\Omega\). If \(|f(z)| \leq M\) for all \(z \in \Omega\) and all \(f \in F\), prove that~\(F\) is equi-continuous on~\(\Omega\).
\def\ProblemHeadingText{Problem 3.}
\item
Let \(f_n \in H(\Omega)\), \(n=1,2,3,\ldots\), be one to one in~\(\Omega\). If \(f_n \to f\) uniformly on compact subsets in~\(\Omega\), prove that~\(f\) is either one to one or constant. Show with examples that both conclusions can occur.
\def\ProblemHeadingText{Problem 4.}
\item
Let \(f \in H(\Omega)\). If~\(f\) has no zeros in~\(\Omega\), prove that \(\ln |f|\) is harmonic in~\(\Omega\).
\def\ProblemHeadingText{Problem 5.}
\item
Evaluate the integral
\[
\int_0^\infty \frac{\ln x}{(x^2+b^2)^2}, \quad b>0.
\]
\def\ProblemHeadingText{Problem 6.}
\item
Let~\(\mathbb{C}\) be the complex plane. Suppose \(f:\mathbb{C}\to\mathbb{C}\) is continuous and~\(f\) is holomorphic except in \([-1,1]\). Prove that~\(f\) is entire.
\def\ProblemHeadingText{Problem 7.}
\item
Prove that a doubly periodic entire function is constant.
\def\ProblemHeadingText{Problem 8.}
\item
Let~\(P\) be a polynomial and~\(C\) the circle \(|z-a|=R\). Evaluate the integral \(\int_C P(z)\,d\overline{z}\).
\def\ProblemHeadingText{Problem 9.}
\item
Let \(f,F \in H(U)\), where \(U=\{z:|z|<1\}\), and let \(R_f=\{f(z):z\in U\}\) and \(R_F=\{F(z):z\in U\}\). Suppose~\(F\) is one to one in~\(U\), \(f(0)=F(0)\), and \(R_f\) is contained in \(R_F\). Prove that there exists an \(\omega \in H(U)\) such that \(f(z)=F(\omega(z))\) and \(|\omega(z)|\leq |z|\). Also show that the equality holds if and only if \(R_f=R_F\).
\def\ProblemHeadingText{Problem 10.}
\item
Suppose \(f \in H(\Pi^+)\), where \(\Pi^+\) is the upper half plane and \(|f|\leq 1\). How large can \(|f'(i)|\) be? Find the extremal functions.
\def\ProblemHeadingText{Problem 11.}
\item
Find all entire functions~\(f\) such that \(|f(z)|=1\) whenever \(|z|=1\).
\end{examproblems}
\end{document}
